\section{Discussion}

\subsection{Principal findings}

The extracted literature supports a multidimensional framework for the digital patient spectrum. The most useful axes are the modeled biological scale, degree of patient specificity, temporal orientation, data-coupling mode, modeling approach, function, validation evidence, and maturity. A definition based only on the phrase \emph{digital twin} is too coarse for healthcare, where a static image-based simulation, an in silico trial cohort, and a continuously updated patient replica may be described with adjacent language.

The central translational gap is not lack of enthusiasm but lack of operational specificity. Reviews repeatedly describe personalized care, prediction, diagnosis, monitoring, and treatment planning, while fidelity metrics, validation datasets, uncertainty analysis, reproducibility, interoperability, governance, and workflow integration are reported unevenly. These gaps make it difficult to distinguish a conceptual prototype, research simulation, decision-support artifact, and operational patient twin.

\subsection{Implications}

Future work should prioritize a reporting and appraisal standard for digital patient reviews and primary models. At minimum, studies should report the physical referent, biological hierarchy, data sources, update frequency, coupling direction, model architecture, validation strategy, uncertainty handling, interoperability standards, governance assumptions, and maturity level. Such a standard would help separate useful analogies from clinically ready infrastructure.

The repository structure demonstrates a practical mechanism for that transparency: protocol and adjudication documents are versioned with the extraction workbook; classification rules are explicit; row-level derived data are inspectable; tables and figures are regenerated by one command; and the manuscript imports generated artifacts directly. This arrangement does not replace substantive team adjudication, but it makes every descriptive count auditable and refreshable as extraction evolves.

\subsection{Limitations}

This synthesis is descriptive and does not estimate pooled effects. The extraction workbook contains free-text, multi-label fields with heterogeneous specificity, and regular-expression classification cannot resolve all semantic ambiguity. Counts are therefore record-level signals rather than independent study effects. Some quality-appraisal and reproducibility fields remain incomplete. Title-level duplicate normalization does not substitute for final bibliographic deduplication. Finally, the PRISMA source counts require arithmetic reconciliation before submission, as documented automatically in the generated quality report.

\subsection{Future work}

The next revision should complete team adjudication, reconcile duplicate records and PRISMA counts, finalize author affiliations and declarations, connect reference-manager records to the LaTeX bibliography, and validate classification rules against a manually reviewed subset. The pipeline can then be rerun without changing the manuscript layout.

\section{Conclusion}

Digital patient research is conceptually active but infrastructurally immature. Progress toward clinically deployable patient digital twins will require clearer taxonomy, stronger and more transparent validation, interoperable data architecture, prospective workflow evidence, and governance models that address privacy, equity, accountability, and reproducibility.
